\section{Introduction}\label{sec:introduction}

Signed adjacency matrices provide a concrete meeting point of spectral graph
theory, switching theory, and discrete magnetic ideas. Switching at vertices
conjugates the signed adjacency matrix by a diagonal sign matrix, so its
spectrum depends on cycle fluxes rather than on an edge gauge. Minimizing the
spectral radius over all signings of a fixed graph is therefore an optimization
problem on a finite flux space. It belongs both to the classical theory of
signed graphs~\cite{Zaslavsky1982} and to the broader spectral-signing problem
initiated by Bilu and Linial~\cite{BiluLinial2006}. The latter line includes the
interlacing-family solution for bipartite Ramanujan lifts and subsequent
explicit near-Ramanujan constructions~\cite{MarcusSpielmanSrivastava2015,MohantyODonnellParedes2022};
those results do not determine the two-sided optimum for a fixed nonbipartite
graph. Extremal-radius questions within signed graph classes form a related,
distinct literature~\cite{BelardoCioabaKoolenWang2019,BelardoBrunettiCiampella2021}.

We study this problem for the four-regular circulant

\begin{equation}
G_n=C_n(1,2),
\end{equation}

whose vertices are the residues modulo \(n\) and whose edges join vertices at
cyclic distances one and two. This is a useful model rather than a universal
one: it is the simplest four-regular circulant with overlapping triangles and
quadrilaterals, its switching space has explicit cycle coordinates, and a
periodic signing admits a finite Floquet reduction. Most importantly for the
present problem, squaring its signed operator makes a local factor \(1+Q_i\)
visible: negative quadrilateral flux cancels a coupling and positive flux
activates it. Thus a global spectral optimization can be related to a local
defect geometry while retaining a complete finite quotient.

Suvagiya recently identified a distinguished family in which the triangle
fluxes alternate~\cite{Suvagiya2026Signed}. For even \(n\), the two members of that
family with negative step-one Hamilton-cycle holonomy have spectral radius

\begin{equation}
\rho_-(n)=2 \sqrt{\cos^{2}(\frac{\pi}{n})+\cos^{2}(\frac{2\pi}{n})}.
\end{equation}

Conjecture 3 of~\cite{Suvagiya2026Signed} states that no signing of \(G_n\) has
smaller spectral radius. That paper verified the assertion for
\(n=8,10,\ldots,18\) and left the general lower bound open. The companion
parity-family framework appears in~\cite{Suvagiya2026Parity}.

Circulant Fourier analysis~\cite{Davis1979} and Floquet theory for periodic
graph operators~\cite{Kuchment1993,KorotyaevSaburova2017} provide the operator
context. Flux-dependent trace formulas on periodic graphs are especially close
in spirit to our closed-walk calculations~\cite{KorotyaevSaburova2023}. We use
``flux phase'' as mathematical terminology suggested by magnetic-operator and
physical precedents, including Lieb's theorem~\cite{Lieb1994}; no physical
ground-state principle or hypothesis from that theorem is invoked here.

The inherited ingredients are the alternating-triangle family, its four
switching/holonomy classes, the displayed formula for \(\rho_-(n)\), and the
conjecture itself~\cite{Suvagiya2026Signed}. The present contribution is not
only an explicit counterexample. We determine the smallest failure, extend it to
an infinite arithmetic family, calculate its exact band edge, identify the
local cancellation mechanism, classify every period-eight phase, derive
general-period obstructions, and complete an exact bounded-period frontier.

The conjecture is false, but its failure is delayed: every admissible
order below 32 satisfies the proposed lower bound. The smallest counterexample is
an order-32 signing whose triangle-flux word is the fourfold repetition of

\begin{equation}
(+,+,-,+,-,-,+,-).
\end{equation}

The same word defines an eight-periodic operator. Its Floquet polynomial can be
computed exactly, and a uniform positivity argument turns the isolated witness
into an infinite family for all multiples of eight at least 32. The resulting
analysis reveals why that word is effective. In the squared operator, entries
with \(Q_i=-1\) cancel and entries with \(Q_i=+1\) become active defects; the
counterexample arranges its two positive defects antipodally in an eight-site
cell. This mechanism, rather than the isolated order-32 comparison, is the
organizing theme of the paper.

\subsection{Main results}

Our first result combines an exact witness with exhaustive finite exclusion.

\begin{theorem}[Smallest Counterexample]
\label{thm:smallest-counterexample}
The conjectured lower bound holds for
every even \(n\) with \(8\leq n\leq 30\). It fails at \(n=32\). Hence 32 is the smallest
counterexample order.

\end{theorem}
The exclusion through 30 is a finite computer-assisted theorem. Its proof
enumerates a complete quotient of the switching classes and certifies every
nonoptimizer by exact rational inequalities. The order-32 witness itself has a
short exact positive-definiteness certificate.

The witness extends periodically.

\begin{theorem}[Infinite Counterexample Family]
\label{thm:infinite-counterexample-family}
For every \(n=8L\) with \(L\geq 4\)
and either Hamilton holonomy, the period-eight signing defined in Section~\ref{sec:periodic-construction}
satisfies

\begin{equation}
\rho(A)^{2} < \frac{1561}{200} < \rho_-(n)^{2}.
\end{equation}

Thus the conjecture fails at every multiple of eight at least 32. We do not
claim that it fails at every even order above 32.

\end{theorem}
The Floquet analysis is sharp at infinite volume.

\begin{theorem}[Exact Period-Eight Edge]
\label{thm:exact-period-eight-edge}
For the target period-eight phase,
the squared infinite-volume spectral radius is

\begin{equation}
\eta=4+\sqrt{10+2\sqrt{5}},
\end{equation}

and the top band reaches \(\eta\) only at Bloch parameter \(z=1\).

\end{theorem}
The local mechanism is described by a complete phase classification. For a
legal eight-periodic flux word \(Q\), let \(D(Q)\) be the set of positive entries
and let \(R(Q)\) denote the squared infinite-volume spectral radius.

\begin{theorem}[Period-Eight Trichotomy]
\label{thm:eight-barrier-trichotomy}
Exactly one of the following occurs:

\begin{equation}
\begin{cases}
R(Q)=8, & D(Q)=\varnothing; \\
R(Q)=\eta<8, & D(Q)=\{j,j+4\}; \\
R(Q)>8, & \text{for every other legal period-eight }Q.
\end{cases}
\end{equation}

In particular the target is the unique period-eight minimizer up to
translation, reflection, and global edge-sign negation, and the all-negative phase
is the unique runner-up.

\end{theorem}
The same squared-operator calculation gives information at every period. Put
\(d=\lvert D(Q)\rvert\), let \(a\) count adjacent positive pairs, and let \(b\) count positive
pairs at cyclic distance two.

\begin{theorem}[General-Period Moment Obstruction]
\label{thm:general-period-moment-obstruction}
For every legal
\(p\)-periodic phase,

\begin{equation}
\begin{aligned}
M_{1}&=4p, \\
M_{2}&=20p+16d, \\
M_{3}&=118p+168d+96a+48b.
\end{aligned}
\end{equation}

If \(R(Q)\leq 8\), then necessarily

\begin{equation}
\begin{aligned}
d&\leq \frac{3p}{4}, \\
40d+96a+48b&\leq 42p.
\end{aligned}
\end{equation}

These are necessary conditions only.

\end{theorem}
Finally we classify a bounded but nontrivial periodic domain.

\begin{theorem}[Low-Period Spectral Classification]
\label{thm:low-period-frontier}
Among the infinite operators \eqref{eq:periodic-operator} whose Hamilton-gauge
triangle word \(\tau\) has least positive period at most 16, the target phase
uniquely minimizes the squared infinite-volume radius \(R\), up to translation,
reflection, global edge-sign negation, and repetition
of the unit cell. Repeated cells are retained only for certificate accounting
and are identified spectrally by zone folding.

\end{theorem}
This theorem is computer-assisted: a complete exact finite classification combines closed-walk compression with exact certificates for the residual competitors. The full partition and certificate accounting appear in Section~\ref{sec:low-period-frontier}.

\subsection{Proof strategy}

The first organizing device is Hamilton gauge. After switching all step-one
edges to \(+1\) except for a single holonomy cut, a signing is encoded by a
periodic word \(\tau\) of triangle fluxes. Its adjacent products

\begin{equation}
Q_i=\tau_i \tau_{i+1}
\end{equation}

are the quadrilateral fluxes. Periodicity permits a Floquet decomposition into
finite Hermitian Laurent matrices \(H_Q(z)\). For the target word, the
characteristic determinant is even in the eigenvalue and reduces to a quartic
\(P(y,c)\) in \(y=\lambda^{2}\) and \(c=z+z^{-1}\). An exact positive-coefficient
expansion at the candidate edge proves the sharp bound and its equality case.

The second device is to square the infinite operator. In \(A_\tau^{2}\), every
odd-displacement coefficient contains a factor \(1+Q_i\). Negative flux
therefore cancels that coupling exactly, while positive flux activates an
amplitude of absolute value two. Constant terms of even Bloch traces count
signed closed walks. If all squared band values were at most eight, their
successive moments would satisfy \(M_{k+1}\leq 8M_k\). A positive excess therefore
forces the spectrum above eight. This one-way implication, combined with
defect geometry, proves the period-eight trichotomy and supplies the general
period obstructions.

\subsection{Scope and public status}

The results distinguish three optimization domains. Theorem~\ref{thm:eight-barrier-trichotomy} concerns
infinite-volume phases of displayed period eight. Theorem~\ref{thm:low-period-frontier} concerns periodic
phases of primitive Hamilton-gauge period at most 16. Theorem~\ref{thm:smallest-counterexample} concerns all
finite signings only at the enumerated orders through 32. None of these results
proves global optimality over all periods, all nonperiodic signings, or every
finite order.

To the best of our knowledge, there is no direct public resolution of Conjecture 3 in
the source paper, the author's listed follow-up materials, or the narrow
direct-result searches recorded for this draft. This is a bounded literature statement, not an absolute priority claim; the search date and scope are recorded in the reproducibility note.

\subsection{Organization}

Section~\ref{sec:preliminaries} fixes switching, flux, Floquet, moment, and equivalence notation.
Section~\ref{sec:smallest-counterexample} proves Theorem~\ref{thm:smallest-counterexample}. Sections~\ref{sec:periodic-construction} and~\ref{sec:period-eight-edge} prove the periodic family and its
sharp edge. Section~\ref{sec:eight-barrier} proves the eight-barrier trichotomy and explains the
chiral symmetry. Section~\ref{sec:general-period} derives the general-period moments and necessary
conditions. Section~\ref{sec:low-period-frontier} proves the bounded low-period frontier. Section~\ref{sec:computational-verification} states
the computer-assisted proof boundary and reproducibility model. Section~\ref{sec:discussion}
collects consequences and open problems. The appendices give the orbit-counting
arguments, exact finite certificates, and computational protocol needed to
audit the computer-assisted parts.
